\documentclass[anonymous]{iopjournal}

% -------------------------
% Packages
% -------------------------
\usepackage[T1]{fontenc}
\usepackage[utf8]{inputenc}
\usepackage{amsmath,amssymb}
\usepackage{graphicx}
\usepackage{booktabs}
\usepackage{array}
\usepackage{tabularx}
\usepackage{threeparttable}
\usepackage{enumitem}
\usepackage{microtype}
\usepackage[section]{placeins}
\usepackage{siunitx}
\usepackage{url}
\usepackage{xcolor} % CHANGED: added for visible TODO markers

% -------------------------
% Formatting controls
% -------------------------
\graphicspath{{figures/}}

\setlength{\parindent}{1.5em}
\setlength{\parskip}{0pt}
\raggedbottom

\setlist[itemize]{topsep=2pt,itemsep=1pt,parsep=0pt,leftmargin=*}
\setlist[enumerate]{topsep=2pt,itemsep=1pt,parsep=0pt,leftmargin=*}

\sisetup{
  detect-all,
  per-mode=symbol,
  range-phrase=--,
  range-units=single
}

% -------------------------
% TODO markers (CHANGED: added). Remove these and all \TODO/\TODOBLOCK uses
% before submission. They render in red so nothing slips through.
% -------------------------
\newcommand{\TODO}[1]{\textcolor{red}{\textbf{[TODO: #1]}}}
\newcommand{\TODOBLOCK}[1]{%
  \par\medskip\noindent
  \fcolorbox{red}{red!8}{%
    \begin{minipage}{0.95\linewidth}
    \vspace{0.3em}
    \textcolor{red}{\textbf{TODO}}\par\smallskip
    #1
    \vspace{0.3em}
    \end{minipage}}%
  \par\medskip
}

% Placeholder-safe figure command.
% This prevents the document from failing if a figure is missing.
\newcommand{\safeinclude}[2][]{%
  \IfFileExists{figures/#2}{%
    \includegraphics[#1]{figures/#2}%
  }{%
    \fbox{%
      \begin{minipage}[c][0.18\textheight][c]{0.88\linewidth}
      \centering
      \vspace{0.5em}
      \textbf{Figure file missing}\\[0.3em]
      Upload \texttt{figures/#2}
      \vspace{0.5em}
      \end{minipage}
    }%
  }%
}

% Compact table column types.
\newcolumntype{Y}{>{\raggedright\arraybackslash}X}
\newcolumntype{L}[1]{>{\raggedright\arraybackslash}p{#1}}

% -------------------------
% Manuscript metadata
% -------------------------
\begin{document}

\articletype{Paper}

\title{A Computational Image-to-Mesh Pipeline for Cardiac Vascular Segmentation and Bioprinting-Oriented STL Generation from CT Angiography}

\author{[Author information removed for double-anonymous review]}

\affil{[Affiliation information removed for double-anonymous review]}

\begin{abstract}
Vascularization remains a central bottleneck in thick engineered tissues, which require reliable oxygen, nutrient, and waste-transport networks. Image-derived patient-specific vascular geometries may support future biofabrication workflows, but translation from clinical scans to fabrication-oriented 3D models remains fragmented across preprocessing, segmentation, mesh generation, repair, and quality control. We present a computational image-to-mesh pipeline linking contrast-enhanced cardiac computed tomography angiography (CTA) to bioprinting-oriented vascular model generation. The workflow includes Phase A, a MONAI/PyTorch volumetric coronary vascular segmentation model, and Phase B, a postprocessing pipeline that converts segmentation outputs into stereolithography (STL) artifacts for downstream review. Phase A uses a three-dimensional Attention U-Net trained on the public ImageCAS coronary CTA dataset with a predefined 700/50/250 train--validation--test split. Preprocessing includes RAS orientation standardization, resampling to \SI{0.6}{mm} isotropic spacing, CT windowing from $-200$ to $700$ HU, normalization to $[0,1]$, binary vessel labels, and foreground-biased $96\times192\times192$ voxel patch sampling. On the completed 250-case held-out test, segmentation achieved mean Dice@0.5 of 0.7953 (95\% CI 0.7896--0.8009), median Dice@0.5 of 0.7985, and zero failed cases; soft Dice, precision, recall, HD95, and runtime were also extracted. Phase B converts probability maps or binary masks into STL outputs using geometric validation, scalar-field generation, isosurface extraction, smoothing, repair, and quality-control reporting. A representative Phase B artifact was watertight and single-component after repair, but failed the available radius-based printability proxy because 39.94\% of sampled radii were below the configured \SI{0.6}{mm} minimum-radius criterion. A representative mesh-smoothing example reduced mean surface roughness from \SI{0.070}{mm} to \SI{0.044}{mm}. These case-level computational demonstrations are not cohort-level physical printability validation. Overall, the results support a computational image-to-mesh workflow with near-0.80 mean Dice@0.5 on a full held-out test set, while topology-aware metrics, cohort-level mesh-quality statistics, slicer-based export validation, physical printing, perfusion, and biological validation remain necessary before stronger translational claims.

\end{abstract}

\vspace{0.8em}
\keywords{3D bioprinting, vascularization, cardiac computed tomography angiography, coronary artery segmentation, 3D U-Net, Attention U-Net, MONAI, PyTorch, STL generation, mesh processing, biofabrication, patient-specific modeling}
\date{}
\maketitle

% CHANGED: master TODO summary (delete before submission)
\TODOBLOCK{This draft is internally consistent and honest, but it is not yet
complete. The items below are blocked on a GPU re-run of inference, because the
final evaluation saved no per-case predictions and recorded
\texttt{compute\_cldice:false}, \texttt{save\_case\_outputs:false}, and
\texttt{run\_phase\_b\_qc:false}. Each is also flagged inline at the point it
belongs.
\begin{enumerate}
  \item \textbf{clDice} across the 250-case test set (\S\ref{subsec:metrics}, Table~\ref{tab:results}).
  \item \textbf{Cohort-level Phase~B mesh QC} (\S\ref{subsec:phaseb-plan}, Table~\ref{tab:phaseb-qc}).
  \item \textbf{Paired segmentation-to-mesh analysis} (\S\ref{subsec:paired}).
  \item Optional, pending a scope decision: a \textbf{non-attention 3D U-Net baseline} under identical splits.
  \item Optional, pending a scope decision: \textbf{external validation} on an independent coronary CTA dataset.
  \item Resolve the \textbf{v14 evaluated-checkpoint provenance} (\S\ref{subsec:training-convergence}): the architecture is proven (strict load, zero key mismatches), but the averaged weights and soup recipe are not reconstructable. An accountable checkpoint is now available (\texttt{checkpoints/best\_dice05.pt}, epoch 79, 23.6M-parameter Attention U-Net, validation Dice@0.5 $=0.789$); the planned resolution is to re-evaluate the 250-case test set from it and replace the headline metrics throughout.
  \item \textbf{Full augmentation parameters} for Appendix~\ref{app:repro}. \emph{Done} --- see Table~\ref{tab:aug}.
  \item Proper formatting of the \textbf{WHO reference} in \texttt{references.bib}. \emph{Done.}
\end{enumerate}}

% =========================================================
\section{Introduction}
% =========================================================

Vascularization remains one of the most persistent constraints in tissue engineering and biofabrication. Thick tissues require connected perfusable networks for oxygen delivery, nutrient transport, and waste removal; without those networks, constructs rapidly become transport-limited and lose viability \cite{murphy2014,lovett2009,rouwkema2016,oconnor2022}. Cardiovascular disease remains the leading global cause of mortality \cite{who,tsao2023}, and engineered cardiac tissues are a long-term motivation for vascularized biofabrication. In cardiac tissue engineering, this challenge is especially important because the functional demands on tissue survival, geometry, and perfusion are unusually strict \cite{murphy2020,noor2019,wu2023}.

Recent advances in extrusion-based bioprinting and freeform reversible embedding of suspended hydrogels have expanded the range of soft, perfusable geometries that can be fabricated \cite{kolesky2014,ozbolat2016,hinton2015}. However, the upstream step of converting patient imaging to fabrication-oriented anatomy remains a bottleneck. Clinical CTA contains spatial information that may be useful for patient-specific vascular templates, yet preparing those data for fabrication typically requires multiple disconnected steps: image preprocessing, segmentation, manual correction, mesh extraction, smoothing, repair, quality control, and export.

Most medical image segmentation studies evaluate models primarily through voxelwise overlap metrics. For biofabrication, that is not sufficient. A vessel mask with an acceptable Dice score may still be unsuitable for manufacturing if it contains discontinuous branches, non-manifold surfaces, geometric misregistration, or wall features that are too thin for stable deposition. Accordingly, segmentation performance and bioprinting-readiness should be treated as related but distinct objectives.

This study presents a computational pipeline that addresses that gap. The goal is not only to segment coronary vascular anatomy from cardiac CTA, but also to convert segmentation outputs into geometrically consistent STL models for downstream fabrication-oriented review. The contribution is deliberately framed as an integration and feasibility study rather than a claim of functional tissue fabrication. First, the pipeline organizes a CTA-to-mesh workflow into reproducible Phase A and Phase B stages. Second, it connects conventional volumetric segmentation with mesh-generation and quality-control outputs. Third, it defines a realistic claim boundary: the current evidence supports computational preparation of print-oriented vascular geometry, not yet physical printing, perfusion, or biological function. By pairing full-cohort segmentation evaluation with representative computational mesh-integrity and fabrication-readiness checks, the study frames image-derived vascular modeling as a coupled segmentation-to-geometry problem rather than a segmentation-only task.

% =========================================================
\section{Background and Related Work}
% =========================================================

Biofabrication has increasingly focused on vascularized constructs because perfusion is fundamental to the viability of larger engineered tissues \cite{groll2016,lovett2009,rouwkema2016,oconnor2022}, with applications spanning cardiac, musculoskeletal, and other regenerative contexts \cite{yang2023}. Patient-specific cardiac bioprinting studies have shown the promise of anatomically tailored tissues and perfusable patches, but they also reveal how difficult it remains to recreate vascular hierarchy at relevant scales \cite{noor2019,wu2023,alonzo2019}. Extrusion-based printing is robust and versatile, whereas FRESH and related embedded approaches improve handling of soft, freeform geometries \cite{ozbolat2016,hinton2015}.

On the imaging side, deep learning has become a major paradigm for biomedical segmentation. U-Net established the encoder--decoder design that has become standard in medical image analysis, and 3D U-Net extended that paradigm to volumetric data \cite{ronneberger2015,cicek2016}. Attention-gated variants improve localization of relevant structures in complex images \cite{oktay2018}. Frameworks such as nnU-Net and MONAI have further advanced reproducible development of volumetric pipelines \cite{isensee2021,cardoso2022}.

For coronary CTA specifically, public datasets have historically been limited, which has made comparisons across studies difficult. ImageCAS helps address that problem by providing a large-scale public benchmark for coronary artery segmentation \cite{zeng2023}. More recent work on ImageCAS has reported coronary segmentation performance using Dice, sensitivity, precision, and HD95, demonstrating the value of reporting multiple complementary metrics rather than Dice alone \cite{huang2025}. Nevertheless, even accurate CTA segmentation is not equivalent to fabrication readiness. Once a segmentation is produced, volumetric masks still require mesh extraction, geometric validation, smoothing, repair, and topology-aware quality checks before they can be treated as fabrication-oriented geometry \cite{lorensen1987,lewiner2003,taubin1995}.

The gap addressed here therefore lies at the interface between medical image analysis and biofabrication. Prior work has advanced segmentation, bioprinting, or mesh processing in isolation. Fewer studies explicitly connect patient-specific CTA segmentation to print-oriented vascular model generation within a single workflow.

% =========================================================
\section{Methods}
% =========================================================

\subsection{Pipeline overview}

The proposed workflow converts cardiac CTA into bioprinting-oriented three-dimensional geometry through two linked stages: (i) automated coronary vascular segmentation and (ii) print-aware postprocessing and STL generation. A schematic overview is shown in figure~\ref{fig:pipeline}. The first stage ingests preprocessed CTA volumes and produces voxelwise vascular predictions. The second stage validates image geometry, standardizes masks or probability maps, reconstructs a surface mesh, improves mesh quality, exports STL models, and generates quality-control outputs for downstream review.

\begin{figure}[h!]
    \centering
    \safeinclude[width=0.86\textwidth]{figure1_pipeline_overview.png}
    \caption{Overview of the CTA-to-STL workflow. The pipeline links preprocessing, volumetric vascular segmentation, probability-map generation, thresholding or scalar-field preparation, mesh extraction, mesh repair, quality-control reporting, and final STL export.}
    \label{fig:pipeline}
\end{figure}

The repository implementation is structured around a Phase A training/inference pipeline and a Phase B mesh-processing pipeline. The full pipeline saves Phase A outputs including \texttt{ct\_preprocessed.nii.gz}, \texttt{seg\_prob.nii.gz}, and optionally \texttt{seg\_mask.nii.gz}, followed by Phase B outputs such as \texttt{vessels\_raw.stl}, \texttt{vessels\_repaired.stl}, \texttt{qc\_report.json}, \texttt{qc\_summary.txt}, \texttt{qc\_summary.csv}, and \texttt{used\_config.yaml}. This organization supports a reproducible image-to-mesh workflow rather than a segmentation-only model.

\subsection{Dataset and preprocessing}

Supervised training was performed on the public ImageCAS dataset, which contains approximately 1000 cardiac CTA cases with coronary artery annotations \cite{zeng2023}. The exported training and evaluation artifacts confirm that 1000 valid image-label samples were discovered. The implementation used an existing \texttt{splits.json} file containing 700 training cases, 50 validation cases, and 250 held-out test cases. Model selection was performed using the 50-case validation partition. The final reported performance in this manuscript comes from the complete 250-case held-out test partition, not from a smaller pilot subset. Split verification found no duplicate cases within splits and no overlap between the training, validation, and held-out test partitions.

Raw DICOM studies were converted to NIfTI prior to training or, during inference, converted when a DICOM directory was supplied \cite{li2016}. Preprocessing was implemented using MONAI transforms. Images and labels were loaded, converted to channel-first format, reoriented to RAS, and resampled to \SI{0.6}{mm} isotropic spacing. For CT inputs, intensities were clipped to $[-200,700]$ HU and linearly scaled to $[0,1]$. Labels were binarized to a single vessel foreground class. Training patches were extracted at $96\times192\times192$ voxels. The logs report a representative raw case shape of $512\times512\times275$ voxels before training-time preprocessing. A representative preprocessing panel showing raw, windowed, and resampled CTA is shown in figure~\ref{fig:preprocessing}.

Data augmentation was applied during training. In the reported A40 run the pipeline comprised spatial padding, foreground/background random cropping (positive-to-negative ratio $3{:}1$), random flipping, random 90-degree rotation, random small-angle spatial rotation, contrast adjustment, intensity scaling, intensity shifting, Gaussian noise, Gaussian smoothing, and a final resize-with-pad-or-crop to the configured patch size. Elastic deformation is implemented in the codebase but was disabled for the reported run because it degraded thin-vessel topology; it was therefore not applied to the evaluated model. The complete per-transform probabilities and magnitude ranges are listed in Appendix~\ref{app:repro} (Table~\ref{tab:aug}). Validation and inference use sliding-window inference with Gaussian blending to reconstruct full-volume predictions.

\begin{table}[!t]
\centering
\caption{Dataset and preprocessing summary.}
\label{tab:dataset}
\begin{threeparttable}
\footnotesize
\begin{tabularx}{0.92\textwidth}{L{0.38\textwidth}Y}
\toprule
\textbf{Item} & \textbf{Value used or reported} \\
\midrule
Primary supervised dataset & ImageCAS \\
Imaging modality & Coronary/cardiac CTA \\
Samples discovered in logs & 1000 image-label cases \\
Confirmed split file & Existing \texttt{splits.json} \\
Training cases & 700 \\
Validation cases & 50 \\
Held-out test cases & 250 \\
Final held-out evaluation & Complete 250-case test partition \\
Failed held-out test cases & 0 \\
Annotation target & Coronary artery / vascular foreground labels \\
Representative raw volume shape & $512\times512\times275$ voxels \\
Preprocessing orientation & RAS \\
Voxel spacing & \SI{0.6}{mm} isotropic \\
CT intensity window & $[-200,700]$ HU \\
Intensity scaling & Linear scaling to $[0,1]$ \\
Training patch size & $96\times192\times192$ voxels \\
Evaluation interpretation & The primary segmentation result is reported on the full held-out 250-case test set after model selection on the 50-case validation set. \\
\bottomrule
\end{tabularx}
\end{threeparttable}
\end{table}

\begin{figure}[!t]
    \centering
    \safeinclude[width=0.72\textwidth]{figure2_preprocessing.png}
    \caption{Representative preprocessing panel showing raw CTA, windowed CTA, and resampled or normalized CTA prepared for model ingestion.}
    \label{fig:preprocessing}
\end{figure}

\FloatBarrier

\subsection{Model architecture}

The submitted segmentation model is defined as a three-dimensional Attention U-Net implemented in MONAI on top of PyTorch \cite{ronneberger2015,cicek2016,oktay2018,cardoso2022,paszke2019}. The model accepts a single-channel volumetric CTA input and outputs a single-channel vessel probability map via sigmoid activation. The configured model uses channel widths of 32, 64, 128, 256, and 512, with strides of 2, 2, 2, and 2.

% CHANGED: lock the evaluated architecture from the exported test console log (resolves the parameter-count ambiguity).
The model evaluated for the final held-out test result in this manuscript is this 23,625,953-parameter configuration. The exported test console log records the loaded checkpoint as a MONAI \texttt{AttentionUnet} with channels $(32,64,128,256,512)$, three residual units per block, and dropout $0.1$, totaling 23,625,953 total and trainable parameters; an independent reconstruction of the same configuration in a separate environment reproduced this exact parameter count. The evaluation load diagnostics additionally record a strict checkpoint load (\texttt{strict=true}) with zero missing and zero unexpected state-dict keys and \texttt{model\_type: attention\_unet}, confirming that the evaluated weights matched this architecture exactly. Other exploratory U-Net-family runs reported 28,650,933 trainable parameters; those runs are treated as development runs and are not used to define the submitted architecture or to report any final result.

\subsection{Training procedure}

Training used supervised learning with patch-based sampling. The exported A40 training configuration records a batch size of 1, gradient accumulation over 2 steps, an ROI size of $96\times192\times192$ voxels, deterministic seed 42, AdamW optimization, gradient clipping with maximum norm 1.0, and a reduce-on-plateau learning-rate schedule in the best exported validation-history run. Training and validation were performed on an NVIDIA A40 GPU. The exported training run used foreground-biased sampling with a 3:1 positive-to-negative patch ratio and two samples per training case; the final held-out evaluation configuration records sliding-window inference with overlap 0.625 and a probability threshold of 0.5.

The project logs show that training labels were highly imbalanced, with 78,936,341 positive vessel voxels and 47,202,404,075 negative voxels in the training set, corresponding to a negative-to-positive voxel ratio of approximately 598:1. This is the true voxel imbalance. Separately, the exported training configuration used a capped or fixed BCE positive weight of 200.0, which should not be interpreted as the empirical class ratio.

% CHANGED: state the evaluated loss plainly.
The evaluated model was trained with a Tversky--Dice composite objective~\cite{salehi2017,sudre2017}, the \texttt{tversky\_fn} configuration recorded in the final-evaluation run name. In this mode the source code combines a Tversky term and a Dice term at full resolution, with an additional downsampled Tversky supervision term. The implementation also supports alternative Dice, BCE, focal, and composite modes, but no formal loss-function ablation is claimed in this study. The loss configuration is therefore reported as part of the implementation rather than as an independently optimized contribution.

% CHANGED: honest note on the v11/v14 provenance split.
The fully documented hyperparameter set above corresponds to the exported \texttt{bioprint\_v11} run, which provides the training- and validation-history records summarized in Section~\ref{subsec:training-convergence}. The checkpoint actually evaluated on the held-out test set was a weight-averaged ``soup'' from a later Tversky--Dice run sharing the same Attention U-Net architecture; its full training-run folder was not retained (Section~\ref{subsec:training-convergence}). \TODO{once the evaluated checkpoint is recovered or re-created (or a re-evaluation from an accountable checkpoint is performed), reconcile this paragraph so the reported test metrics and the documented training run refer to the same checkpoint.}

\begin{table}[!t]
\centering
\caption{Training configuration supported by the implementation and exported logs.}
\label{tab:training}
\begin{threeparttable}
\footnotesize
\begin{tabularx}{0.92\textwidth}{L{0.42\textwidth}Y}
\toprule
\textbf{Setting} & \textbf{Value} \\
\midrule
Frameworks & MONAI / PyTorch \\
Submitted architecture & MONAI \texttt{AttentionUnet} / 3D Attention U-Net family \\
Attention U-Net parameter count & 23,625,953 trainable parameters \\
Residual units / dropout & 3 / 0.1 \\ % CHANGED: added from console log
Input channels & 1 \\
Output channels & 1 \\
Network channels & 32, 64, 128, 256, 512 \\
Strides & 2, 2, 2, 2 \\
Patch size & $96\times192\times192$ voxels \\
Batch size / gradient accumulation & 1 / 2 steps \\
Training / validation / test cases & 700 / 50 / 250 \\
Positive vessel voxels in training labels & 78,936,341 \\
Negative background voxels in training labels & 47,202,404,075 \\
Negative-to-positive voxel ratio & Approximately 598:1 \\
Capped/fixed positive weight in exported v11 run & 200.0 \\
Foreground-biased sampling in exported v11 run & 3:1 positive-to-negative patch ratio; 2 samples per case \\
Evaluated loss objective & Tversky--Dice composite (\texttt{tversky\_fn}) \\ % CHANGED
Optimizer & AdamW \\
Learning-rate schedule & Reduce-on-plateau in exported v11 run \\
Gradient clipping & Maximum norm 1.0 \\
Final training hardware & NVIDIA A40 GPU \\
Checkpoint selection & Validation Dice@0.5; final test evaluated once on the held-out partition \\
\bottomrule
\end{tabularx}
\end{threeparttable}
\end{table}

\FloatBarrier

\subsection{Phase A inference workflow}

The full inference pipeline loads a trained checkpoint, reconstructs the model, applies the same inference preprocessing, and performs sliding-window inference with Gaussian blending. When the input is a DICOM directory, the pipeline first converts it to NIfTI. The model output is converted to probabilities using sigmoid activation. The Phase A output directory stores the preprocessed CT volume, vessel probability map, and, when a threshold is provided, an optional binary mask:
\begin{itemize}
    \item \texttt{ct\_preprocessed.nii.gz}
    \item \texttt{seg\_prob.nii.gz}
    \item \texttt{seg\_mask.nii.gz}
\end{itemize}
These outputs provide the bridge between the segmentation stage and the downstream mesh-generation stage.

\subsection{Bioprinting-oriented postprocessing and mesh generation}

The second stage is what makes the workflow biofabrication-relevant rather than purely diagnostic. Instead of terminating at voxel masks, the pipeline transforms segmentations into surface models that can be inspected, repaired, and exported for downstream fabrication-oriented workflows. The contribution of Phase B is not novelty in marching cubes or smoothing alone; rather, it is the integration of segmentation output, metadata validation, mesh extraction, repair, and QC reporting into a single image-to-STL pipeline.

\subsubsection{Auxiliary anatomical segmentation}

TotalSegmentator v1.5.5 was used only as an auxiliary tool for cardiac-region localization or anatomical context during preprocessing and downstream spatial organization \cite{wasserthal2023}. It was not used to generate coronary artery labels, was not treated as a microvascular reconstruction method, and was not used as evidence of capillary-scale vascular modeling. This distinction is important because the supervised vascular output of this study comes from the ImageCAS-trained coronary segmentation model, not from TotalSegmentator.

\subsubsection{Geometric validation and mask standardization}

Before mesh generation, segmentation masks are checked against the reference CTA for spacing, origin, and orientation consistency. Let $\mathbf{s}_{\mathrm{mask}}$ and $\mathbf{s}_{\mathrm{ref}}$ denote the voxel-spacing vectors of the mask and reference image. The spacing validation criterion described in the draft can be written as
\begin{equation}
\left\lVert \mathbf{s}_{\mathrm{mask}} - \mathbf{s}_{\mathrm{ref}} \right\rVert_{\infty}
\leq \tau_s,
\qquad
\tau_s = 10^{-3}\ \mathrm{mm}.
\label{eq:spacing}
\end{equation}
Direction matrices are checked with an elementwise tolerance of $10^{-3}$. Multi-label masks are then binarized to foreground/background representations and assigned reference metadata to preserve voxel-to-world correspondence.

\subsubsection{Thickness-aware shell generation}

For structures intended to provide printable bulk geometry, the draft describes morphology-based shell generation with a target wall thickness of \SI{3}{mm}. With target thickness $t_{\mathrm{mm}}$ and voxel spacings $(s_x,s_y,s_z)$, the dilation radius in voxel units is
\begin{equation}
r_{\mathrm{vox}}
=
\max
\left(
1,
\left\lceil
\frac{t_{\mathrm{mm}}}{\min(s_x,s_y,s_z)}
\right\rceil
\right).
\label{eq:dilation}
\end{equation}
A spherical dilation is applied to the original mask, and the printable shell is obtained by set subtraction between the dilated and original objects. Because cohort-level wall-thickness results are not reported here, this step is described as a fabrication-oriented processing component rather than validated physical printability.

\subsubsection{Scalar-field generation and isosurface extraction}

To reduce staircase artifacts, binary masks are converted to continuous scalar fields before surface reconstruction. For discrete masks, Gaussian smoothing is applied after binarization. For probabilistic inputs, floating-point values are retained directly. The scalar field is optionally supersampled and padded prior to isosurface extraction. Surface meshes are then reconstructed using marching cubes, specifically the Lewiner variant with topological guarantees \cite{lorensen1987,lewiner2003}. The resulting triangulation is mapped back into physical coordinates using the scan spacing and origin.

\subsubsection{Volume-preserving mesh optimization and quality control}

Meshes are subsequently smoothed with Taubin's alternating low-pass procedure to reduce high-frequency roughness while minimizing global shrinkage \cite{taubin1995}. The pipeline is designed to generate raw and repaired STL files and quality-control artifacts, including JSON, text, and CSV summaries. These outputs are intended to support review of watertightness, manifoldness, connectedness, repair status, and related fabrication-readiness checks. In this manuscript, mesh-processing results are reported as representative computational demonstrations rather than cohort-level validation of printability.

\begin{figure}[!t]
    \centering
    \safeinclude[width=0.70\textwidth]{figure3_taubin_smoothing.png}
    \caption{Representative Taubin-smoothing example. In the illustrated case, maximum roughness decreased from \SI{0.237}{mm} to \SI{0.120}{mm}, and mean roughness decreased from \SI{0.070}{mm} to \SI{0.044}{mm}. This single-case result demonstrates the intended smoothing effect but is not interpreted as cohort-level mesh-quality validation.}
    \label{fig:taubin}
\end{figure}

\FloatBarrier


\subsection{Computational fabrication-readiness validation}

To evaluate whether segmentation-derived STL outputs were suitable for downstream fabrication-oriented review, the Phase B workflow included a computational fabrication-readiness validation stage. This stage assessed mesh integrity, spatial fidelity, and basic print-preparation constraints without claiming physical bioprinting, perfusion, or biological function. Mesh integrity was evaluated using watertightness, boundary-edge status, manifold topology, and connectedness checks. A mesh was considered geometrically valid for downstream review when no open boundary edges were detected and when edge connectivity did not indicate non-manifold topology. Wall-thickness compliance was evaluated against a minimum structural threshold of \SI{3}{mm} for the generated shell or printable geometry, where applicable. Because this threshold is a fabrication-oriented structural constraint rather than a biological coronary-wall measurement, it was interpreted as a print-preparation criterion rather than a claim of anatomical wall-thickness realism.

Slicability was assessed using a multi-plane cross-sectional continuity check. In this procedure, the mesh was intersected with 50 cross-sectional planes, and each slice was inspected computationally for continuous closed contours. Mesh-quality stability was further assessed using edge-length distributions and comparison against the target edge-length range. A radius-based printability proxy was additionally computed, comparing sampled local surface radii against a configurable minimum-radius criterion (here \SI{0.6}{mm}). Finally, spatial fidelity was evaluated by verifying DICOM alignment preservation, voxel-spacing preservation, and center-of-mass consistency between the image-derived segmentation and exported STL geometry. These checks were designed to determine whether the generated STL artifacts were geometrically coherent and suitable for downstream slicing or phantom-printing experiments. They do not establish physical print success, lumen patency, perfusion, cell viability, endothelialization, or biological function.

\begin{table}[!t]
\centering
\caption{Computational fabrication-readiness checks performed on representative Phase B STL outputs. These checks evaluate geometric validity and print-preparation suitability but do not constitute physical printing, perfusion, or biological validation.}
\label{tab:fabrication-readiness-checks}
\begin{threeparttable}
\footnotesize
\begin{tabularx}{0.96\textwidth}{L{0.31\textwidth}Y}
\toprule
\textbf{Check} & \textbf{Interpretation} \\
\midrule
Watertightness & Evaluates whether the exported STL surface is closed without open boundary edges. \\
Boundary-edge count & Detects holes or leaks that may prevent slicing or downstream geometric processing. \\
Manifold topology & Evaluates whether mesh edges have valid face connectivity and avoid non-manifold defects. \\
Wall-thickness compliance & Checks whether the generated shell or printable geometry satisfies the selected minimum structural threshold. \\
Minimum-radius (printability) proxy & Compares sampled local surface radii against a configurable minimum-radius criterion to flag features below typical deposition scale. \\ % CHANGED: added the proxy used in results
Slicability & Uses multi-plane cross-sectional inspection to assess whether the geometry produces continuous contours during slicing. \\
Edge-length distribution & Quantifies triangle-size stability and identifies excessive mesh irregularity after reconstruction or smoothing. \\
DICOM alignment preservation & Confirms that exported geometry remains spatially aligned with the source image coordinate system. \\
Voxel-spacing preservation & Verifies that mesh conversion does not introduce spacing-related geometric distortion. \\
Center-of-mass validation & Checks whether the exported mesh remains anatomically aligned with the source segmentation. \\
\bottomrule
\end{tabularx}
\end{threeparttable}
\end{table}

\FloatBarrier

\subsection{Phase B cohort-level mesh-quality assessment plan}
\label{subsec:phaseb-plan}

Beyond the representative computational checks reported here, the geometry stage should ultimately be reported through cohort-level mesh-quality summaries aggregated across the evaluation partition. Case-level QC artifacts already produced by the workflow, including \texttt{qc\_report.json}, \texttt{qc\_summary.txt}, and \texttt{qc\_summary.csv}, should be consolidated into cohort tables so that mesh usability can be assessed across cases rather than only by example. Recommended cohort-level aggregation files are \texttt{outputs/phase\_b\_mesh\_qc/per\_case\_mesh\_qc.csv}, \texttt{outputs/phase\_b\_mesh\_qc/summary\_mesh\_qc.csv}, and \texttt{outputs/phase\_b\_mesh\_qc/summary\_mesh\_qc.json}. If the same folder is not used in the final implementation, the manuscript should still preserve this one-file-per-purpose structure.

\TODOBLOCK{This subsection and Table~\ref{tab:phaseb-qc} are still \emph{planned} rather than \emph{reported}. Running Phase~B across all 250 held-out cases (on the locked predictions) will populate every metric below with per-case values and cohort summaries. This is blocked only by the GPU re-run noted at the top of the manuscript. Once complete, convert this from a plan into a Results subsection and fill Table~\ref{tab:phaseb-qc}.}

\begin{table}[!t]
\centering
\caption{Planned cohort-level Phase B mesh-quality metrics for revised submission. These metrics are defined here to make clear which fabrication-oriented properties will be reported once cohort aggregation is completed.}
\label{tab:phaseb-qc}
\begin{threeparttable}
\footnotesize
\begin{tabularx}{0.96\textwidth}{L{0.30\textwidth}Y}
\toprule
\textbf{Metric} & \textbf{Interpretation and planned reporting} \\
\midrule
Connected-component count & Number of disconnected mesh components after STL generation and repair; reported per case and summarized across the cohort. \TODO{fill} \\
Watertightness & Boolean or categorical indicator of whether the STL surface is closed and suitable for downstream geometric processing. \TODO{fill (\% watertight)} \\
Non-manifold edge count & Number of topologically invalid edges remaining after mesh extraction or repair. \TODO{fill} \\
Repair success & Whether the repair stage resolves holes, self-intersections, or non-manifold artifacts sufficiently for export. \TODO{fill (success rate)} \\
Surface roughness & Pre- and post-smoothing roughness summaries, including mean and maximum roughness where available. \TODO{fill} \\
Smoothing-induced volume change & Relative change in mesh volume after smoothing to document shape preservation versus shrinkage. \TODO{fill} \\
Minimum-radius proxy & Fraction of sampled radii below the selected minimum-radius criterion, summarized across the cohort. \TODO{fill} \\ % CHANGED
Wall-thickness compliance & Fraction of the mesh satisfying the selected minimum fabrication-oriented wall-thickness threshold. \TODO{fill} \\
Slicer compatibility & Whether the exported STL imports successfully into the target slicing or print-preparation software without fatal geometry errors. \TODO{fill} \\
QC artifact paths & Case-level source files retained as \texttt{qc\_report.json} and \texttt{qc\_summary.csv}; cohort aggregation retained as \texttt{outputs/phase\_b\_mesh\_qc/per\_case\_mesh\_qc.csv}. \\
\bottomrule
\end{tabularx}
\end{threeparttable}
\end{table}

\subsection{Evaluation metrics}
\label{subsec:metrics}

Evaluation explicitly separates segmentation accuracy from fabrication readiness. For segmentation, the completed held-out test evaluation computed soft Dice, Dice@0.5, precision@0.5, recall@0.5, 95th-percentile Hausdorff distance (HD95), foreground voxel counts, false-positive voxels, false-negative voxels, predicted foreground voxels, and per-case runtime. Dice@0.5 was retained as the primary metric because it was the predefined thresholded metric used for checkpoint selection and reporting. Confidence intervals in the exported summary file are reported as 95\% intervals over the 250 per-case values.

The validation logs contain threshold sweeps across probability thresholds from 0.1 to 0.9. These validation sweeps are useful for diagnosing calibration and recall/precision tradeoffs. The final held-out test evaluation exported only threshold 0.5 metrics. Therefore, the manuscript reports Dice@0.5 as the primary final test result and does not tune the primary threshold on the held-out test set.

For tubular vascular anatomy, topology-sensitive measures such as clDice are especially relevant because they emphasize centerline overlap and connectedness more directly than volumetric overlap alone \cite{shit2021}. The completed held-out test evaluation did not compute clDice (the evaluation configuration recorded \texttt{compute\_cldice:false} and did not save per-case predictions), so this manuscript does not yet report a topology metric. \TODO{compute clDice@0.5 across the 250-case test set on the locked-model predictions and report mean/median/95\% CI here and in Table~\ref{tab:results}. This is the single most important missing metric for a paper whose thesis is that voxel overlap is insufficient for tubular geometry; it is blocked only by the GPU re-run.} Topology-aware evaluation is treated as an essential validation step before stronger biofabrication-readiness claims can be made.

For the geometry stage, the relevant outputs are connected, watertight, printable meshes rather than masks alone. Accordingly, the most relevant mesh-quality quantities include connected-component count, non-manifold edge count, watertightness, repair success, minimum wall thickness, minimum-radius proxy, edge-length distributions, surface-roughness summaries, and slicer compatibility. Representative computational fabrication-readiness checks are reported for a Phase B output in Section~\ref{subsec:mesh-results}, but they are not reported as cohort-level outcomes and are not treated as physical print validation in the present manuscript.

\begin{table}[!t]
\centering
\caption{Evaluation metrics and quality-control outputs.}
\label{tab:evaluation}
\begin{threeparttable}
\footnotesize
\begin{tabularx}{0.92\textwidth}{L{0.35\textwidth}Y}
\toprule
\textbf{Category} & \textbf{Metrics or outputs} \\
\midrule
Primary held-out segmentation metric & Dice@0.5 on the 250-case test set \\
Additional held-out segmentation metrics & Soft Dice, precision@0.5, recall@0.5, HD95@0.5, FP/FN voxels, predicted and ground-truth foreground voxels, runtime \\
Validation calibration diagnostics & Dice at thresholds 0.1--0.9, probability diagnostics, and voxel counts above thresholds where available in validation logs \\
Tubular topology metric & clDice identified as necessary; not computed in the completed held-out test evaluation \TODO{compute and report} \\
Postprocessing cleanup & Connected-component filtering supported in implementation \\
Phase A output files & \texttt{ct\_preprocessed.nii.gz}, \texttt{seg\_prob.nii.gz}, optional \texttt{seg\_mask.nii.gz} \\
Phase B output files & \texttt{vessels\_raw.stl}, \texttt{vessels\_repaired.stl}, \texttt{qc\_report.json}, \texttt{qc\_summary.txt}, \texttt{qc\_summary.csv}, \texttt{used\_config.yaml} \\
Mesh-quality validation targets & Watertightness, non-manifold edge count, connected-component count, repair success rate, wall-thickness compliance, minimum-radius proxy, and surface roughness distributions \\
\bottomrule
\end{tabularx}
\end{threeparttable}
\end{table}

\subsection{Comparative and paired analyses}
\label{subsec:paired}

The completed results reported here focus on the selected Attention U-Net model and the full held-out 250-case test evaluation. A direct baseline comparison against a non-attention 3D U-Net or nnU-Net-style reference was not completed within the present manuscript, so no superiority claim is made. Future comparative work should preserve the same \texttt{splits.json} partition, preprocessing, ROI size, sliding-window inference settings, checkpoint-selection rule, and held-out evaluation procedure so that architectural differences are not confounded by evaluation drift.

A paired segmentation-to-mesh analysis is also needed so that fabrication-oriented usability can be evaluated at the case level rather than inferred from overlap metrics alone. For each test case, segmentation metrics such as soft Dice, threshold-specific Dice, precision, recall, HD95, clDice, false-positive voxels, false-negative voxels, and predicted-versus-ground-truth foreground voxel counts should be linked to mesh outcomes including connected-component count, watertightness, non-manifold edge count, repair success, surface roughness, minimum-radius proxy, wall-thickness compliance, and slicer compatibility.

\TODOBLOCK{The paired analysis is the most novel contribution and is currently described as ``needed'' rather than performed. Once the locked-model per-case segmentation metrics (with clDice) and the cohort Phase~B mesh QC both exist, join them by case ID and report Spearman/Pearson correlations (with $p$-values) identifying which segmentation metrics predict mesh usability, plus scatterplots for the strongest relationships. This converts the ``coupled segmentation-to-geometry'' framing in the Introduction from a claim into a result. Blocked only by the GPU re-run.}

% =========================================================
\section{Results}
% =========================================================

\subsection{Training convergence and model selection}
\label{subsec:training-convergence}

The strongest exported validation-history record is the A40 run \texttt{bioprint\_v11\_a40\_roi96\_posneg3to1\_lr2e5}. In this run, validation Dice@0.5 improved from 0.7146 at epoch 62 to a best exported value of 0.7542 at epoch 73, with a corresponding validation soft Dice of 0.7080. The same run used the verified 700/50/250 split, Attention U-Net architecture with 23,625,953 trainable parameters, $96\times192\times192$ voxel patches, foreground-biased patch sampling, and capped positive weighting for class imbalance. Figure~\ref{fig:training-curves} summarizes the exported training and validation trajectory.

% CHANGED: document the v14 soup as unrecoverable, with explicit remediation.
The final held-out evaluation configuration records a checkpoint-soup path, \texttt{/tmp/train\_runs/bioprint\_v14\_tversky\_fn/checkpoints/soup\_ep75-80.pt}, evaluated on the test partition. The naming indicates a weight-averaged ensemble of epoch-75 through epoch-80 checkpoints from a Tversky--Dice run sharing the architecture in Section~3.3. This checkpoint resided on ephemeral cloud storage (\texttt{/tmp}) and was not retained; a search of the cloud run directory, the local repository, and the download archive did not locate the soup file, its epoch-75--80 ingredient checkpoints, or a soup-construction script. The evaluation record nonetheless documents what was run: the configuration specifies \texttt{model\_type: attention\_unet} and the load diagnostics report a strict checkpoint load with empty missing- and unexpected-key lists, so the evaluation provably used the 23,625,953-parameter Attention U-Net of Section~3.3 even though the averaged weights themselves are no longer retained. The manuscript therefore reports the held-out test result directly from the saved per-case and summary evaluation files, which remain available and auditable, and avoids unsupported claims about the detailed construction of the checkpoint soup.
\TODOBLOCK{The evaluated \emph{architecture} is now established (the load diagnostics show a strict load with zero key mismatches, see above); what remains unreconstructable is the averaged \emph{weights} and the soup recipe. Resolve before final deposition, in order of preference: (a) recover the soup or its epoch-75--80 ingredient checkpoints from any surviving backup of the A40/cloud environment and document the averaging recipe in Appendix~\ref{app:provenance}; (b) if only ingredients survive, re-create the soup (uniform average of epochs 75--80) and confirm it reproduces the reported metrics; (c) if nothing survives, re-evaluate the 250-case test set once from an accountable checkpoint (e.g. the best-by-validation v11 checkpoint, or a documented v11 soup) and replace the headline metrics throughout. In all three cases, enable \texttt{save\_case\_outputs}, \texttt{compute\_cldice}, and \texttt{run\_phase\_b\_qc} so the re-run also unblocks clDice and the cohort Phase~B QC.}

\begin{figure}[!t]
    \centering
    \safeinclude[width=0.86\textwidth]{figure4_training_curves.png}
    \caption{Training and validation curves for the exported A40 validation-history run. The curves show train Dice, validation Dice@0.5, validation soft Dice, training loss, and validation loss across epochs.}
    \label{fig:training-curves}
\end{figure}

\subsection{Full 250-case held-out test performance}

The completed held-out test evaluation included all 250 test cases in \texttt{splits.json} and recorded zero failed cases. The primary metric, Dice@0.5, had a mean of 0.7953 (95\% CI 0.7896--0.8009), median 0.7985, standard deviation 0.0457, IQR 0.0599, and range 0.6143--0.8715. The mean soft Dice was 0.7886 (95\% CI 0.7830--0.7942). Precision@0.5 and recall@0.5 were balanced, with mean precision 0.7961 (95\% CI 0.7880--0.8041; median 0.8081) and mean recall 0.8008 (95\% CI 0.7924--0.8093; median 0.8097). Mean HD95@0.5 was \SI{4.92}{mm} (median \SI{3.52}{mm}); the HD95 distribution was strongly right-skewed (SD \SI{5.66}{mm}, maximum \SI{46.09}{mm}), indicating that although most cases achieved close boundary agreement, a small number exhibited large boundary errors consistent with distal-branch or thin-vessel failures. Predicted and ground-truth foreground volumes were closely matched on average (mean 31,450 versus 31,629 voxels), and false-positive and false-negative voxels were of comparable magnitude (means 6,280 and 6,459), consistent with the balanced precision and recall. Mean per-case inference runtime was \SI{5.39}{s} (median \SI{5.64}{s}, range 1.69--10.69\,s); because Phase B mesh QC was not run during this evaluation, this runtime reflects Phase A inference only. The held-out partition comprises ImageCAS cases 751--1000. Full per-metric distributions are summarized in table~\ref{tab:full-metrics}, and the per-case Dice@0.5 distribution is shown in figure~\ref{fig:test-dice-distribution}.

\begin{table}[!t]
\centering
\caption{Full per-metric statistics over the 250-case held-out test set, read directly from \texttt{summary\_metrics.json}. Confidence intervals are 95\% intervals over the 250 per-case values.}
\label{tab:full-metrics}
\begin{threeparttable}
\footnotesize
\setlength{\tabcolsep}{5pt}
\begin{tabularx}{0.97\textwidth}{L{0.24\textwidth} l l l Y}
\toprule
\textbf{Metric} & \textbf{Mean (95\% CI)} & \textbf{Median} & \textbf{SD} & \textbf{Range (min--max)} \\
\midrule
Dice@0.5 & 0.7953 (0.7896--0.8009) & 0.7985 & 0.0457 & 0.6143--0.8715 \\
Soft Dice & 0.7886 (0.7830--0.7942) & 0.7927 & 0.0455 & 0.6035--0.8649 \\
Precision@0.5 & 0.7961 (0.7880--0.8041) & 0.8081 & 0.0650 & 0.5385--0.9322 \\
Recall@0.5 & 0.8008 (0.7924--0.8093) & 0.8097 & 0.0684 & 0.5898--0.9402 \\
HD95@0.5 (mm) & 4.92 (4.22--5.62) & 3.52 & 5.66 & 0.50--46.09 \\
Inference runtime (s/case) & 5.39 (5.17--5.61) & 5.64 & 1.77 & 1.69--10.69 \\
\midrule
GT foreground (voxels) & 31,629 (30,423--32,835) & 30,485 & 9,728 & 8,783--64,439 \\
Predicted (voxels) & 31,450 (30,412--32,489) & 30,623 & 8,376 & 9,266--57,976 \\
False positive (voxels) & 6,280 (5,998--6,563) & 5,919 & 2,278 & 1,936--14,870 \\
False negative (voxels) & 6,459 (6,035--6,884) & 5,802 & 3,424 & 968--21,247 \\
\bottomrule
\end{tabularx}
\begin{tablenotes}\footnotesize
\item clDice is intentionally absent: it was not computed in this evaluation (\texttt{compute\_cldice:false}) and requires the re-run noted in the text.
\end{tablenotes}
\end{threeparttable}
\end{table}

\begin{figure}[!t]
    \centering
    \safeinclude[width=0.74\textwidth]{figure7_full_test_dice_distribution.png}
    \caption{Distribution of Dice@0.5 across the full held-out 250-case ImageCAS test set. Vertical reference lines denote the mean and median.}
    \label{fig:test-dice-distribution}
\end{figure}

\subsection{Threshold sensitivity and calibration}

Validation logs contained threshold sweeps from 0.1 to 0.9. In the exported v11 validation run, the best sweep value occurred near high probability thresholds, with the best exported validation sweep Dice of 0.7918 at threshold 0.866. This suggests that the validation predictions were relatively conservative and that threshold calibration can materially affect apparent Dice. However, the final held-out test evaluation was run at threshold 0.5 only. To avoid test-set threshold tuning, Dice@0.5 remains the primary held-out test metric, and threshold sensitivity is treated as a validation-set diagnostic rather than a redefinition of the final test endpoint.

\subsection{Case-level variability}

Per-case Dice@0.5 varied across the held-out cohort, with a minimum of 0.6143, maximum of 0.8715, and IQR of 0.0599. The worst cases therefore remained substantially below the cohort mean even though the median approached 0.80. This variability is expected in coronary CTA segmentation because distal branches, reduced contrast, motion artifacts, and annotation uncertainty can disproportionately affect thin vessels. Ranked per-case Dice and the relationship between ground-truth foreground volume and Dice are shown in figures~\ref{fig:ranked-dice} and~\ref{fig:volume-dice}.

\begin{figure}[!t]
    \centering
    \safeinclude[width=0.78\textwidth]{figure8_per_case_dice_ranked.png}
    \caption{Per-case Dice@0.5 values ranked across the full held-out 250-case test set. The plot illustrates cohort-level robustness and residual failure modes.}
    \label{fig:ranked-dice}
\end{figure}

\begin{figure}[!t]
    \centering
    \safeinclude[width=0.70\textwidth]{figure9_foreground_volume_vs_dice.png}
    \caption{Ground-truth vessel foreground volume versus Dice@0.5 across the held-out test set. This diagnostic plot is included to show whether segmentation performance is associated with vessel-label volume.}
    \label{fig:volume-dice}
\end{figure}

\subsection{Benchmark context}

The final held-out Dice@0.5 value is placed in context with public ImageCAS-related literature in table~\ref{tab:benchmark}. The comparison is not a direct head-to-head benchmark because published methods may use different splits, preprocessing, model configurations, postprocessing, and metric definitions. Nevertheless, the completed 250-case result places the computational pipeline in a plausible ImageCAS performance range while preserving the claim boundary that this study is an image-to-mesh feasibility and integration study, not a state-of-the-art segmentation benchmark.

\begin{table}[!t]
\centering
\caption{Contextual comparison with ImageCAS-related coronary CTA segmentation work. Values are included for benchmark context only and are not treated as directly comparable unless dataset split, preprocessing, and metric definitions match.}
\label{tab:benchmark}
\begin{threeparttable}
\footnotesize
\setlength{\tabcolsep}{4pt}

\renewcommand{\arraystretch}{1.18}
\begin{tabularx}{0.96\textwidth}{L{0.26\textwidth}L{0.30\textwidth}Y}
\toprule
\textbf{Study or method} & \textbf{Evaluation context} & \textbf{Reported performance and interpretation} \\
\midrule
ImageCAS benchmark study \cite{zeng2023} &
Public ImageCAS coronary CTA benchmark &
Introduced the large-scale ImageCAS dataset and benchmark framework for coronary artery segmentation. \\
\addlinespace[0.25em]
Anatomy-guided ImageCAS method \cite{huang2025} &
ImageCAS with 7:1:2 train/validation/test split &
Reported Dice 0.8082, sensitivity 0.7946, precision 0.8471, and HD95 \SI{9.77}{mm}. This provides recent benchmark context with multiple metrics. \\
\addlinespace[0.25em]
This work &
ImageCAS-derived 700/50/250 split; full 250-case held-out test set &
Mean Dice@0.5 0.7953 (95\% CI 0.7896--0.8009), median 0.7985, zero failed cases. This supports near-0.80 computational segmentation performance but is not presented as a direct state-of-the-art comparison. \\
\bottomrule
\end{tabularx}
\end{threeparttable}
\end{table}

\subsection{Mesh generation and computational fabrication-readiness checks}
\label{subsec:mesh-results}

A representative postprocessing example for mesh smoothing is shown in figure~\ref{fig:taubin}. In that example, maximum roughness decreases from \SI{0.237}{mm} to \SI{0.120}{mm} and mean roughness decreases from \SI{0.070}{mm} to \SI{0.044}{mm}, corresponding to a 37.1\% improvement in mean roughness after Taubin smoothing. Because cohort-level mesh-quality distributions are not reported, this result is interpreted only as an illustrative demonstration of the mesh-processing stage.

The automated full-pipeline script supports Phase A inference followed by Phase B mesh generation. The completed 250-case evaluation quantified segmentation performance and per-case Phase A inference runtime (mean \SI{5.39}{s} per case; table~\ref{tab:full-metrics}) but did not run cohort-level Phase B mesh QC, so Phase B runtime is not yet available. \TODO{report cohort Phase~B QC and Phase~B runtime after the GPU re-run, alongside the Phase~A runtime already recorded.}

% CHANGED: honest verified case-720 result, reframed as scale mismatch.
A representative computational fabrication-readiness panel was generated from the Phase B outputs for case~720 (figure~\ref{fig:fabrication-readiness-qc}). The input to this Phase~B demonstration was a binary mask thresholded at 0.4 (the Phase~B mesh-generation threshold, set independently of the 0.5 threshold used for the reported segmentation metrics), yielding a single mask component. The reconstructed geometry was watertight both before and after repair and consisted of a single connected mesh component, with a reconstructed surface area of 7854.84\,mm$^2$ and volume of 7071.36\,mm$^3$. The overall computational QC status for this case was reported as a failure, driven entirely by the radius-based printability proxy: the minimum mesh radius was \SI{0.355}{mm} and 39.94\% of sampled radii fell below the configured \SI{0.6}{mm} minimum-radius criterion. Because coronary arteries, and distal branches in particular, are intrinsically sub-millimeter, this outcome is better interpreted as evidence that coronary vascular geometry falls below a generic extrusion-printability threshold than as a loss of mesh integrity; the same artifact remained watertight and single-component after repair. This single-case result is interpreted as a representative computational demonstration and not as cohort-level fabrication validation. \TODO{after the cohort Phase~B run, report the distribution of the below-radius fraction across all 250 cases; if most cases fall below the 0.6\,mm proxy, state that the proxy threshold is mismatched to coronary caliber and consider reporting it as a radius distribution rather than a binary pass/fail.}

Representative vessel-network and shell STL renderings from the geometry stage are shown in figures~\ref{fig:vessel-stl} and~\ref{fig:heart-stl}.

\begin{table}[!t]
\centering
\caption{Representative Phase B computational fabrication-readiness results for the evaluated STL output (case~720). Values are read directly from the exported \texttt{qc\_report.json} and \texttt{qc\_summary.csv}. The pass/fail wording is a case-level computational geometry check and not physical printing, perfusion, or biological validation.}
\label{tab:phaseb-representative-results}
\begin{threeparttable}
\footnotesize
\begin{tabularx}{0.92\textwidth}{L{0.36\textwidth}Y}
\toprule
\textbf{Metric or check} & \textbf{Representative observed result (case 720)} \\
\midrule
Segmentation input & Binary mask, threshold 0.4; single mask component \\
Raw mesh watertight & True \\
Repaired mesh watertight & True \\
Mesh components after repair & 1 (single connected component) \\
Reconstructed surface area & 7854.84\,mm$^2$ \\
Reconstructed volume & 7071.36\,mm$^3$ \\
Minimum mesh radius & \SI{0.355}{mm} (criterion \SI{0.6}{mm}) \\
Below-radius fraction (printability proxy) & 39.94\% of sampled radii below \SI{0.6}{mm} \\
Overall computational QC status & Fail (driven by the radius proxy; see text) \\
Manifold topology, slicability, wall-thickness, spatial-fidelity & \TODO{re-run Phase~B with QC enabled to record these for case~720 and across the cohort; the exported panel for this case reports the rows above only} \\
\bottomrule
\end{tabularx}
\end{threeparttable}
\end{table}

\begin{figure}[htbp]
    \centering
    \safeinclude[width=0.82\textwidth]{figure10_fabrication_readiness_qc.png}
    \caption{Representative computational fabrication-readiness panel for a Phase B STL output (case~720). The panel reports raw and repaired watertightness, connected-component count, reconstructed surface area and volume, minimum mesh radius, and the fraction of sampled radii below the \SI{0.6}{mm} minimum-radius printability proxy. For this case the geometry was watertight and single-component after repair, but 39.94\% of radii fell below the proxy (minimum radius \SI{0.355}{mm}), yielding an overall computational QC failure that reflects the sub-millimeter caliber of coronary vessels relative to a generic threshold rather than a loss of mesh integrity. This is a single-case computational demonstration and not evidence of physical printing, perfusion, or biological function.}
    \label{fig:fabrication-readiness-qc}
\end{figure}

\begin{table}[!t]
\centering
\caption{Numerical results and run details available from exported project materials.}
\label{tab:results}
\begin{threeparttable}
\footnotesize
\begin{tabularx}{0.92\textwidth}{L{0.42\textwidth}Y}
\toprule
\textbf{Metric or run detail} & \textbf{Reported or derived value} \\
\midrule
Total samples discovered & 1000 \\
Training / validation / held-out test split & 700 / 50 / 250 \\
Final held-out test cases evaluated & 250 \\
Failed held-out test cases & 0 \\
Primary test metric & Dice@0.5 \\
Test Dice@0.5, mean (95\% CI) & 0.7953 (0.7896--0.8009) \\
Test Dice@0.5, median / SD / IQR & 0.7985 / 0.0457 / 0.0599 \\
Test Dice@0.5, min--max & 0.6143--0.8715 \\
Test soft Dice, mean (95\% CI) & 0.7886 (0.7830--0.7942) \\
Test precision@0.5, mean (95\% CI) & 0.7961 (0.7880--0.8041) \\
Test recall@0.5, mean (95\% CI) & 0.8008 (0.7924--0.8093) \\
Test HD95@0.5, mean / median & \SI{4.92}{mm} / \SI{3.52}{mm} \\
Test runtime, mean / median & \SI{5.3883}{s} / \SI{5.6350}{s} per case \\
clDice (250-case test) & \TODO{compute and report mean/median/95\% CI} \\ % CHANGED
Submitted Attention U-Net parameter count & 23,625,953 trainable parameters \\
Training class imbalance & 78,936,341 positive voxels vs. 47,202,404,075 negative voxels \\
Corrected negative-to-positive ratio & Approximately 598:1 \\
Representative max roughness change & \SI{0.237}{mm} to \SI{0.120}{mm} \\
Representative mean roughness change & \SI{0.070}{mm} to \SI{0.044}{mm} \\
Representative mean roughness improvement & 37.1\% \\
% CHANGED: verified case-720 fabrication-readiness row
Representative Phase B result (case 720) & Watertight and single-component after repair; surface area 7854.84\,mm$^2$; volume 7071.36\,mm$^3$; minimum radius \SI{0.355}{mm}; 39.94\% of radii below the \SI{0.6}{mm} proxy; overall computational QC fail \\
Final test files & \texttt{outputs/final\_test\_250/per\_case\_metrics.csv}, \texttt{summary\_metrics.csv}, \texttt{summary\_metrics.json} \\
Not reported as final metrics & clDice, external-dataset performance, cohort-level Phase B mesh QC statistics, slicer-export success logs, physical printing, perfusion, and biological validation \\
\bottomrule
\end{tabularx}
\end{threeparttable}
\end{table}

\begin{figure}[htbp]
    \centering
    \safeinclude[width=0.62\textwidth]{figure5_vessel_stl_views.png}
    \caption{Representative vessel STL renderings from multiple viewpoints. This figure illustrates the intended output of the mesh-generation stage and the preservation of gross vascular geometry after STL export.}
    \label{fig:vessel-stl}
\end{figure}

\begin{figure}[htbp]
    \centering
    \safeinclude[width=0.62\textwidth]{figure6_heart_stl_views.png}
    \caption{Representative heart or shell STL renderings from multiple viewpoints. This panel provides a qualitative visualization of the downstream geometry workflow and is not presented as evidence of physical printing, perfusion, or biological validation.}
    \label{fig:heart-stl}
\end{figure}

% CHANGED: flag figure legibility/anonymity check for the STL renderings.
\TODO{confirm figures~\ref{fig:vessel-stl} and~\ref{fig:heart-stl} are legible at print size and that no file paths, usernames, or other identifying text are baked into the axis annotations.}

\FloatBarrier

% =========================================================
\section{Discussion}
% =========================================================

\subsection{Principal findings}

This study is best understood as a computational bridge between coronary CTA segmentation and fabrication-oriented mesh generation. The completed full held-out 250-case ImageCAS test evaluation strengthens the evidence for the segmentation stage compared with the earlier pilot framing. The primary Dice@0.5 result was 0.7953 with a 95\% CI of 0.7896--0.8009, median 0.7985, and zero failed cases. Because the confidence interval narrowly spans 0.80, the most accurate wording is near-0.80 segmentation performance rather than a categorical claim that the system exceeds 0.80.

The main implication is that image-to-fabrication pipelines should not terminate at segmentation metrics. For coronary and vascular structures, distal branch dropout, false bridges, and local discontinuities may have limited effect on Dice while substantially affecting downstream geometry. This is especially important for biofabrication, where a disconnected or non-manifold vascular model may fail as a perfusable blueprint even if its voxel overlap is numerically acceptable. The representative Phase~B result makes a related point in the opposite direction: a mesh can be geometrically clean (watertight and single-component) yet still fall outside a generic printability criterion because coronary calibre is intrinsically sub-millimeter. Both observations argue that segmentation accuracy and fabrication readiness must be reported as distinct objectives. The present pipeline begins to address that gap by generating STL artifacts and QC summaries, but the current manuscript stops short of claiming validated printability.

The comparison with ImageCAS-related literature also clarifies the status of the result. The completed held-out test result is plausible relative to public coronary CTA segmentation work and is now supported by per-case statistics, precision, recall, HD95, and runtime. However, the absence of external-dataset validation, clDice, cohort-level mesh QC, and direct baseline training under identical conditions means the method should not be framed as state-of-the-art. Instead, the strongest supported claim is that the pipeline is feasible, reproducible in structure, and aligned with a biofabrication-motivated downstream use case.

\subsection{Biofabrication relevance and claim boundary}

The Phase B contribution is integration rather than invention of new meshing algorithms. Marching cubes and Taubin smoothing are established tools \cite{lorensen1987,lewiner2003,taubin1995}. Their value here comes from being placed within a CTA-to-STL workflow that preserves spatial metadata, accepts segmentation probability maps, produces repaired STL outputs, and records QC artifacts. For a biofabrication audience, this integration matters because the practical barrier is often not any single algorithm, but the reliable movement from medical image volume to usable geometry.

For biofabrication, patient-specific imaging-derived vascular geometry is useful in at least three ways. First, it can serve as a template for perfusable channel design or sacrificial-path planning in engineered tissues \cite{kolesky2014,hinton2015}. Second, it can support fabrication of cardiovascular phantoms or preoperative models where geometric realism matters. Third, it can standardize the upstream computational workflow needed before investigators attempt physical printing. The postprocessing choices described in this study are therefore not incidental. Geometric validation reduces the chance of spatial distortion. Scalar-field reconstruction reduces voxel aliasing that would otherwise degrade surface quality. Volume-preserving smoothing suppresses staircase artifacts without obvious global shrinkage. Mesh repair and QC reporting provide a route toward checking whether segmentation-derived geometry is suitable for downstream fabrication workflows.

The representative computational fabrication-readiness checks strengthen the biofabrication relevance of the pipeline by evaluating whether segmentation-derived STL outputs satisfy basic geometric requirements for downstream print preparation. Watertightness, boundary-edge absence, manifold topology, slicability, edge-length stability, wall-thickness compliance, the minimum-radius proxy, and spatial alignment preservation are necessary intermediate criteria between voxelwise segmentation and physical fabrication. The case-720 result shows how these criteria behave for true coronary geometry: integrity criteria can be satisfied while a deposition-scale criterion is not, which is itself informative about the fabrication challenge for sub-millimeter vasculature. However, these criteria should be interpreted as computational readiness checks rather than evidence of successful bioprinting. Physical print trials, slicer-export validation, lumen patency testing, and perfusion validation are still required before the workflow can be described as a validated vascular bioprinting pipeline.

At the same time, the paper maintains a clear claim boundary. The present results show computational preparation of STL geometry from CTA-derived vascular segmentations. They do not show successful cell-laden printing, perfusion under flow, endothelialization, long-term construct viability, implantation, or clinical readiness. This distinction is central to the interpretation of the work.

\subsection{Limitations}

Several limitations should be stated explicitly. First, the supervised dataset labels coronary arteries rather than the entire cardiac vascular hierarchy. Therefore, the current system should be described as a coronary or macrovascular CTA-to-STL pipeline, not as a full microvascular reconstruction method. Second, the current study is computational: no physical bioprinting experiment, mechanical test, perfusion assay, or biological validation is reported.

Third, while the full 250-case held-out test evaluation is now complete, it remains a single-dataset evaluation. External validation on independent coronary CTA data is still needed. Fourth, Dice@0.5, precision, recall, and HD95 do not prove vessel topology, distal branch continuity, lumen preservation, or fabrication readiness. clDice and centerline-based branch metrics were not computed in the completed held-out test evaluation. Fifth, the Phase B fabrication-readiness evidence in this manuscript is a single representative case rather than a cohort-level result; it does not measure slicer-generated toolpaths, print fidelity, construct mechanics, lumen patency, flow behavior, cell survival, or endothelialization. Sixth, the minimum-radius printability proxy used here (\SI{0.6}{mm}) is a generic deposition-scale threshold and is not calibrated to coronary caliber; the case-720 ``failure'' should be read in that light, and cohort-level radius distributions are needed to interpret it properly. Seventh, cohort-level mesh-quality statistics are not yet available in the current manuscript. Eighth, the final test evaluation records a v14 checkpoint-soup path, but the checkpoint and its construction metadata resided on ephemeral storage and were not recoverable from the cloud run directory, local repository, or download archive; this provenance gap is documented in Section~\ref{subsec:training-convergence} and must be resolved before final deposition.

\subsection{Future validation priorities}

The most important next validation step is topology-aware evaluation of the existing held-out predictions, including clDice, centerline distance, branch detection accuracy, distal-vessel recall, and connected-component behavior \cite{shit2021}. These metrics are necessary because coronary arteries are tubular structures where branch continuity can be more relevant to downstream fabrication than voxel overlap alone.

On the geometry side, the revised study should report full mesh-quality distributions across the evaluation cohort. Important quantities include connected-component count, non-manifold edge count, watertightness, repair success, surface roughness, minimum-radius proxy, wall-thickness compliance, and slicer compatibility. The most useful next analysis would connect segmentation errors to mesh-generation failures so that the field can understand which image-analysis metrics actually predict fabrication-oriented usability.

For translational relevance, the most realistic next fabrication-oriented validations are cohort-level STL repair assessment, formal slicer compatibility testing, slicer-export or toolpath generation logs, minimum-feature or wall-thickness checks, and, if feasible, non-biological phantom printing. Hydrogel or sacrificial-ink fabrication, lumen patency measurements, and ultimately perfusion or cell-seeding studies would further strengthen the translational argument, but those extensions should follow only after the current CTA pipeline is fully documented and quantitatively validated. A particularly useful analysis would link segmentation metrics such as Dice, recall, HD95, and foreground volume to mesh outcomes such as watertightness, non-manifold edge count, slicability, and repair success so that downstream fabrication failures can be related back to image-analysis errors.

% =========================================================
\section{Conclusion}
% =========================================================

This manuscript presents a computational framework that links cardiac CTA to print-oriented vascular model generation. The completed full held-out 250-case ImageCAS test evaluation supports a cautious feasibility claim: volumetric deep learning can automate coronary vascular extraction with near-0.80 Dice@0.5 performance and no failed test cases, and the Phase B workflow provides a structured route from segmentation output to repaired STL geometry and QC artifacts.

The primary held-out test result was a mean Dice@0.5 of 0.7953 (95\% CI 0.7896--0.8009), with median 0.7985 across 250 cases. A representative Phase~B artifact was watertight and single-component after repair, while a radius-based printability proxy flagged its sub-millimeter features, illustrating that geometric integrity and deposition-scale printability are distinct properties for coronary geometry. The contribution is therefore best understood not as a claim of printed organ function, but as a patient-specific image-to-mesh workflow addressing a real bottleneck at the interface of AI, medical imaging, and biofabrication. Topology-aware metrics, cohort-level mesh QC, the paired segmentation-to-mesh analysis, baseline comparisons, external validation, formal slicer validation, physical printing, perfusion testing, and biological validation will be required before stronger translational biofabrication claims can be made.

\section*{Data and code availability}

This study used the public ImageCAS coronary CTA dataset for supervised training and evaluation. The dataset is available from the original ImageCAS source cited in this manuscript. The completed final held-out test artifacts are retained in \texttt{outputs/final\_test\_250/}, including \texttt{per\_case\_metrics.csv}, \texttt{summary\_metrics.csv}, \texttt{summary\_metrics.json}, \texttt{used\_config.yaml}, \texttt{eval\_command.txt}, \texttt{progress.json}, and \texttt{test\_console.log}. These files support verification of the reported 250-case held-out test results. Training and validation artifacts are retained under \texttt{outputs/train\_runs/}, and the representative Phase~B geometry-integrity artifacts for case~720 are retained with the exported Phase~B outputs.

% CHANGED: describe what IS anonymized; honest provenance caveat.
The implementation was developed in Python using PyTorch, MONAI, NumPy, SciPy, NiBabel, pandas, and related medical image-processing libraries. The manuscript text, tables, and figures have been anonymized for double-anonymous review, and an anonymized code and derived-results archive will be provided during review or upon reasonable request, subject to dataset-use restrictions. Raw local training and evaluation logs may still contain author-identifying metadata, usernames, local file paths, institutional names, email addresses, and hard-coded personal directory names; these are not claimed to be anonymized and will be scrubbed prior to any public deposition. \TODO{before camera-ready/deposition, run an anonymity scrub over the manuscript, all figures, and the code/results archive (grep for names, institution, emails, usernames, absolute home paths) and confirm zero hits; then update this paragraph to state that the deposited archive is fully anonymized.}

\section*{Declarations}

\textbf{Ethics.} This study used de-identified public imaging data and did not involve new human-subject recruitment, intervention, clinical deployment, animal experimentation, or collection of private patient data.

\textbf{Competing interests.} The authors declare no competing interests.

\textbf{Acknowledgements.} Author names, institutional identifiers, and funding details have been removed for double-anonymous review and can be restored after peer review if requested by the journal.

% =========================================================
% CHANGED: appendices added (delete if you prefer to match the original exactly).
% =========================================================
\appendix

\section{Reproducibility details}
\label{app:repro}

\subsection*{A.1\quad Training-time augmentation parameters}

Table~\ref{tab:aug} lists every training-time transform applied to the evaluated model, in execution order, with its sampling probability and parameter values as configured in the training code (\texttt{get\_transforms}, \texttt{mode="train"}). All intensity transforms operate on images already scaled to $[0,1]$ (Section~\ref{subsec:dataset}); spatial transforms are applied jointly to image and label with nearest-neighbour interpolation on the label. Deterministic preprocessing common to training, validation, and inference (RAS reorientation, resampling to \SI{0.6}{mm} isotropic spacing, $[-200,700]$ HU windowing, linear scaling to $[0,1]$, and label binarization) is described in Section~\ref{subsec:dataset} and is not repeated here.

\begin{table}[!t]
\centering
\caption{Training-time data-augmentation transforms (MONAI dictionary transforms), in execution order, as used for the evaluated model. Probabilities and parameters are read directly from the training configuration. Elastic deformation is implemented but was disabled for the reported run.}
\label{tab:aug}
\small
\begin{tabular}{@{}lll@{}}
\hline
Transform (MONAI) & Probability & Parameters \\
\hline
\texttt{SpatialPadd} & 1.0 & target $96\times192\times192$ \\
\texttt{RandCropByPosNegLabeld} & --- & pos:neg $=3{:}1$, num\_samples $=2$, allow\_smaller, image\_threshold $=0$ \\
\texttt{RandFlipd} & 0.5 & spatial\_axis $=[0,1,2]$ \\
\texttt{RandRotate90d} & 0.5 & max\_k $=3$ \\
\texttt{RandRotated} & 0.1 & range\_x/y/z $=\pm0.2$ rad ($\approx\pm11.5^\circ$); mode (bilinear, nearest) \\
\texttt{Rand3DElasticd} & disabled & implemented as $p=0.2$, sigma\_range $(4,6)$, magnitude\_range $(50,150)$; not applied \\
\texttt{RandAdjustContrastd} & 0.4 & gamma $\in[0.8,1.2]$ \\
\texttt{RandScaleIntensityd} & 0.5 & factors $=(0.85,1.15)$ \\
\texttt{RandShiftIntensityd} & 0.3 & offsets $=0.1$ \\
\texttt{RandGaussianNoised} & 0.2 & mean $=0.0$, std $=0.01$ \\
\texttt{RandGaussianSmoothd} & 0.3 & sigma\_x/y/z $=(0.5,1.5)$ \\
\texttt{ResizeWithPadOrCropd} & 1.0 & target $96\times192\times192$ \\
\hline
\end{tabular}
\end{table}

\subsection*{A.2\quad Per-case results table}

Per-case held-out test results are provided as \texttt{outputs/final\_test\_250/per\_case\_metrics.csv} (one row per test case). Column legend: \texttt{case\_id}; \texttt{image}, \texttt{label} (source NIfTI paths); \texttt{status}; \texttt{soft\_dice}; \texttt{gt\_voxels} (ground-truth foreground voxels); \texttt{runtime\_seconds}; and, at the reporting threshold $0.5$, \texttt{dice@0.5}, \texttt{precision@0.5}, \texttt{recall@0.5}, \texttt{fp\_voxels@0.5}, \texttt{fn\_voxels@0.5}, \texttt{pred\_voxels@0.5}, \texttt{hd95@0.5} (mm), and \texttt{cldice@0.5}. When Phase~B QC is enabled, per-case mesh-quality columns (connected-component count, watertightness, non-manifold edge count, repair status, surface-roughness, minimum-radius proxy, and wall-thickness compliance) are appended to the same row and mirrored in \texttt{outputs/phase\_b\_mesh\_qc/per\_case\_mesh\_qc.csv}.

\section{Model-provenance record}
\label{app:provenance}

The held-out test metrics were produced by the evaluation command recorded in \texttt{outputs/final\_test\_250/eval\_command.txt}, which loaded the checkpoint \texttt{soup\_ep75-80.pt} and ran sliding-window inference (ROI $96\times192\times192$, overlap 0.625, threshold 0.5) over the 250-case test partition, reporting 0 failed cases and a mean Dice@0.5 of 0.7953. The exported test console log records the loaded model as a MONAI Attention U-Net with channels $(32,64,128,256,512)$, three residual units, dropout 0.1, and 23,625,953 trainable parameters; an independent reconstruction reproduced this parameter count. The evaluation load diagnostics report \texttt{model\_type: attention\_unet} and a strict state-dict load with empty missing- and unexpected-key lists, confirming the evaluated weights matched this architecture. The averaged checkpoint and its epoch-75--80 ingredients were not retained (Section~\ref{subsec:training-convergence}). An accountable, fully retained checkpoint is, however, available for re-evaluation: \texttt{checkpoints/best\_dice05.pt} is an epoch-79 Attention U-Net with 23{,}631{,}875 parameters (matching the architecture above), recorded validation Dice@0.5 $=0.7889$ and soft Dice $=0.7811$, and complete optimizer, scheduler, configuration, and git-commit metadata. The planned resolution is to re-evaluate the 250-case held-out test set from this checkpoint (with \texttt{save\_case\_outputs}, \texttt{compute\_cldice}, and \texttt{run\_phase\_b\_qc} enabled) and replace the headline metrics throughout; the orchestration for this run is provided as \texttt{run\_full\_revision.sh} in the code archive. \TODO{after the re-evaluation, replace the soup path above with \texttt{best\_dice05.pt}, update all per-case metrics/figures, and confirm the strict-load diagnostics for the new run.}

\newpage
% WHO reference (references.bib, key `who`) now has a full @misc entry
% (organization, title, fact-sheet note, year, url, urldate).
\bibliographystyle{iopart-num}
\bibliography{references}

\end{document}